\documentclass[11pt]{article}

\usepackage{amsmath,amssymb,amsthm,mathtools}
\usepackage[margin=1in]{geometry}
\usepackage{microtype}

\newtheorem{theorem}{Theorem}[section]
\newtheorem{lemma}[theorem]{Lemma}
\newtheorem{corollary}[theorem]{Corollary}
\newtheorem{remark}[theorem]{Remark}

\newcommand{\dd}{\,\mathrm d}

\title{Odd Cycles with One Pendant Leaf:\\
Fractional Degree Bias and H\"older Compression}
\author{}
\date{}

\begin{document}

\maketitle

\begin{abstract}
Let $m\geq3$ be odd and let $C_m^+$ be obtained from the odd cycle $C_m$
by attaching one pendant leaf at one cycle vertex.  We prove that every
graphon $W$ of edge density $p\geq1/2$ satisfies
\[
 t(C_m^+,W)
 \geq p\bigl(p^m-p(1-p)^{m-1}\bigr)
 =\Phi_{C_m^+}(p).
\]
The proof averages the leaf root over the cycle, applies the arithmetic--
geometric mean inequality, and changes the underlying probability measure
by the fractional degree weight $d^{1/m}$.  A sharp three-factor H\"older
inequality controls the new edge density, while a one-variable monotonicity
lemma transfers the accepted odd-cycle bound without loss.  The case $m=5$
classifies the formerly open NetworkX Atlas graph $104$.
\end{abstract}


%%%%%%%%%%%%%%%%%%%%%%%%%%%%%%%%%%%%%%%%%%%%%%%%%%%%%%%%%%%%%%%%%%%%%%%%%%%%%%%
\section{Definitions and the exact target}
%%%%%%%%%%%%%%%%%%%%%%%%%%%%%%%%%%%%%%%%%%%%%%%%%%%%%%%%%%%%%%%%%%%%%%%%%%%%%%%

Let $(\Omega,\mu)$ be an arbitrary probability space and let
$W\colon\Omega^2\to[0,1]$ be symmetric and measurable.  Write
\[
 p=\int_{\Omega^2}W(x,y)\dd\mu(x)\dd\mu(y),
 \qquad
 d(x)=\int_\Omega W(x,y)\dd\mu(y).
\]

For odd $m\geq3$, label the vertices of $C_m$ cyclically by
$0,\ldots,m-1$.  Let $C_m^+$ be the graph obtained by adding one new leaf
adjacent to vertex $0$.  The chromatic polynomial of an odd cycle is
\[
 \chi_{C_m}(x)=(x-1)^m-(x-1).
\]
The leaf has $x-1$ choices after its neighbour is coloured, so
\begin{equation}
 \chi_{C_m^+}(x)
 =(x-1)^{m+1}-(x-1)^2,
 \qquad \chi(C_m^+)=3.
 \label{eq:chromatic}
\end{equation}
Consequently, with $q=1-p$,
\begin{align}
 \Phi_{C_m^+}(p)
 &=p\bigl(p^m-pq^{m-1}\bigr)
 \notag\\
 &=p^{m+1}-p^2(1-p)^{m-1}.
 \label{eq:target}
\end{align}
This formula is polynomial and gives the continued value one at $p=1$.

We use the accepted odd-cycle graphon theorem in the following exact form:
for every odd $m\geq3$, every graphon $K$ on every probability space, and
its edge density $s$,
\begin{equation}
 t(C_m,K)\geq\varphi_m(s),
 \qquad
 \varphi_m(s):=s^m-s(1-s)^{m-1}.
 \label{eq:odd-cycle-input}
\end{equation}
This is proved in \texttt{notes/paper\_new\_region2\_v3.tex}.  We do not
assume any rooted strengthening of it.


%%%%%%%%%%%%%%%%%%%%%%%%%%%%%%%%%%%%%%%%%%%%%%%%%%%%%%%%%%%%%%%%%%%%%%%%%%%%%%%
\section{Fractional edge H\"older inequality}
%%%%%%%%%%%%%%%%%%%%%%%%%%%%%%%%%%%%%%%%%%%%%%%%%%%%%%%%%%%%%%%%%%%%%%%%%%%%%%%

The next lemma is useful independently of cycles.

\begin{lemma}[Fractionally degree-weighted edge]\label{lem:fractional-edge}
Let $0<\alpha\leq1$ and put
\[
 N_\alpha
 =\int_{\Omega^2}W(x,y)d(x)^\alpha d(y)^\alpha
 \dd\mu(x)\dd\mu(y).
\]
Then
\[
 N_\alpha\geq p^{1+2\alpha}.
\]
In particular, for $\alpha=1/m$,
\[
 N_{1/m}\geq p^{1+2/m}.
\]
\end{lemma}

\begin{proof}
On the set where $d(x)=0$, one has $W(x,y)=0$ for almost every $y$.
Define $W(x,y)/d(x)$ to be zero there.  Put
\[
 \theta=\frac1{1+2\alpha},
 \qquad
 \eta=\frac{\alpha}{1+2\alpha}.
\]
Then $\theta+2\eta=1$, and almost everywhere on the support of $W$,
\[
 W(x,y)
 =\bigl(W(x,y)d(x)^\alpha d(y)^\alpha\bigr)^\theta
  \left(\frac{W(x,y)}{d(x)}\right)^\eta
  \left(\frac{W(x,y)}{d(y)}\right)^\eta.
\]
Weighted H\"older gives
\[
 p\leq N_\alpha^\theta
 \left(\int\frac{W(x,y)}{d(x)}\dd\mu^2\right)^\eta
 \left(\int\frac{W(x,y)}{d(y)}\dd\mu^2\right)^\eta.
\]
Each quotient integral is at most one: integrating first in the other
variable gives one on $\{d>0\}$ and zero on $\{d=0\}$.  Hence
$p\leq N_\alpha^\theta$, which is the claimed inequality.
\end{proof}


%%%%%%%%%%%%%%%%%%%%%%%%%%%%%%%%%%%%%%%%%%%%%%%%%%%%%%%%%%%%%%%%%%%%%%%%%%%%%%%
\section{The scalar transfer}
%%%%%%%%%%%%%%%%%%%%%%%%%%%%%%%%%%%%%%%%%%%%%%%%%%%%%%%%%%%%%%%%%%%%%%%%%%%%%%%

\begin{lemma}[Odd-cycle rescaling]\label{lem:scalar}
Let $m\geq3$, $p\in[1/2,1]$, and $0<z\leq p$.  Set
\[
 s_0=p\left(\frac pz\right)^{2/m}.
\]
If $s_0\leq1$, then
\[
 z\,\varphi_m(s_0)\geq p\,\varphi_m(p).
\]
\end{lemma}

\begin{proof}
For $s\in(0,1]$, put
\[
 \psi_m(s)=\frac{\varphi_m(s)}{s^{m/2}}
 =s^{m/2}-s^{1-m/2}(1-s)^{m-1}.
\]
Differentiation gives
\begin{align*}
 \psi_m'(s)={}&\frac m2s^{m/2-1}
 +(m-1)s^{1-m/2}(1-s)^{m-2}\\
 &+\left(\frac m2-1\right)s^{-m/2}(1-s)^{m-1}>0.
\end{align*}
Thus $\psi_m$ is strictly increasing.  Write $r=z/p\leq1$.  Then
$s_0=pr^{-2/m}\geq p$ and $r=(p/s_0)^{m/2}$.  Therefore
\[
 z\varphi_m(s_0)
 =p^{1+m/2}\psi_m(s_0)
 \geq p^{1+m/2}\psi_m(p)
 =p\varphi_m(p).
\]
\end{proof}

We shall also use that $\varphi_m$ is increasing on $[1/2,1]$.  Indeed,
\[
 \varphi_m'(s)
 =ms^{m-1}-(1-s)^{m-1}
 +(m-1)s(1-s)^{m-2}>0
\]
there, since $ms^{m-1}\geq m(1-s)^{m-1}$.


%%%%%%%%%%%%%%%%%%%%%%%%%%%%%%%%%%%%%%%%%%%%%%%%%%%%%%%%%%%%%%%%%%%%%%%%%%%%%%%
\section{The one-leaf odd-cycle theorem}
%%%%%%%%%%%%%%%%%%%%%%%%%%%%%%%%%%%%%%%%%%%%%%%%%%%%%%%%%%%%%%%%%%%%%%%%%%%%%%%

\begin{theorem}[One leaf on an odd cycle]\label{thm:main}
For every odd integer $m\geq3$ and every graphon $W$ of edge density
$p\in[1/2,1]$,
\[
 \boxed{
 t(C_m^+,W)
 \geq p\bigl(p^m-p(1-p)^{m-1}\bigr)
 =\Phi_{C_m^+}(p).
 }
\]
\end{theorem}

\begin{proof}
Let
\[
 F(x_0,\ldots,x_{m-1})
 =\prod_{i=0}^{m-1}W(x_i,x_{i+1}),
 \qquad x_m=x_0.
\]
Integrating the pendant leaf first gives
\[
 t(C_m^+,W)=\int F\,d(x_0)\dd\mu^m.
\]
Cyclic relabelling shows that the same density results when the degree
factor is placed at any $x_i$.  Averaging these $m$ equal integrals and
applying arithmetic--geometric mean pointwise gives
\begin{equation}
 t(C_m^+,W)
 \geq\int F\prod_{i=0}^{m-1}d(x_i)^{1/m}\dd\mu^m.
 \label{eq:agm}
\end{equation}

Put
\[
 M=\int_\Omega d(x)^{1/m}\dd\mu(x).
\]
Since $p\geq1/2$, one has $M>0$.  Define the probability measure
\[
 \dd\nu(x)=\frac{d(x)^{1/m}}M\dd\mu(x).
\]
The right side of \eqref{eq:agm} is
\begin{equation}
 M^m t(C_m,W;\nu),
 \label{eq:measure-change}
\end{equation}
where the notation means that all vertex variables are integrated using
$\nu$.  The edge density of $W$ on this new probability space is
\[
 s=\frac{N_{1/m}}{M^2}.
\]

Set $z=M^m$.  Concavity of $r\mapsto r^{1/m}$ gives
\[
 M\leq p^{1/m},
 \qquad 0<z\leq p.
\]
Lemma~\ref{lem:fractional-edge} gives
\[
 s\geq\frac{p^{1+2/m}}{M^2}
 =p\left(\frac pz\right)^{2/m}=:s_0.
\]
Since $W\leq1$ and $\nu$ is a probability measure, $s\leq1$; hence also
$s_0\leq1$.  In particular $s\geq s_0\geq p\geq1/2$.

Apply the accepted odd-cycle theorem \eqref{eq:odd-cycle-input} on
$(\Omega,\nu)$, monotonicity of $\varphi_m$, and
Lemma~\ref{lem:scalar}:
\begin{align*}
 t(C_m^+,W)
 &\geq M^m t(C_m,W;\nu)\\
 &\geq z\varphi_m(s)
 \geq z\varphi_m(s_0)
 \geq p\varphi_m(p).
\end{align*}
The last expression is exactly \eqref{eq:target}.
\end{proof}

At $p=1/2$ the target is zero, while every step remains valid.  At $p=1$,
$W=1$ almost everywhere and both sides equal one.  No division by
$2p-1$ or by a cycle density occurs.

Every balanced complete $k$-partite graphon, $k\geq2$, is regular of degree
$p=1-1/k$.  Then $M=p^{1/m}$, $z=p$, $s=p$, and equality holds in every
step precisely when it holds in the accepted odd-cycle theorem.  Thus the
new theorem has the expected chromatic equality examples.


%%%%%%%%%%%%%%%%%%%%%%%%%%%%%%%%%%%%%%%%%%%%%%%%%%%%%%%%%%%%%%%%%%%%%%%%%%%%%%%
\section{Catalogue consequence and formalization boundary}
%%%%%%%%%%%%%%%%%%%%%%%%%%%%%%%%%%%%%%%%%%%%%%%%%%%%%%%%%%%%%%%%%%%%%%%%%%%%%%%

\begin{corollary}[Atlas 104]\label{cor:atlas104}
NetworkX Atlas graph $104$, with graph6 code \texttt{Ehd?}, is positive.  It
has order six, size six, chromatic number three, and
\[
 \chi_{F_{104}}(x)=(x-1)^6-(x-1)^2.
\]
Every graphon of edge density $p\in[1/2,1]$ satisfies
\[
 t(F_{104},W)
 \geq p^6-p^2(1-p)^4
 =p^2(2p-1)(2p^2-2p+1).
\]
\end{corollary}

\begin{proof}
Its NetworkX edge list is
\[
 01,\ 04,\ 12,\ 15,\ 23,\ 34.
\]
Vertices $0,1,2,3,4$ induce the cycle
$0-1-2-3-4-0$, and vertex $5$ is a single leaf adjacent to vertex $1$.
Thus it is $C_5^+$, and Theorem~\ref{thm:main} applies.
\end{proof}

The proof uses ordinary, non-injective homomorphism densities.  All
integrands are bounded and measurable, so Tonelli--Fubini applies on an
arbitrary probability space.  The quotient $W/d$ is explicitly defined to
be zero on $\{d=0\}$, where $W$ vanishes in the relevant sections almost
everywhere.  No finite-step reduction, injective-copy interpretation,
regularity, or minimum-degree assumption is used.

For formal verification, the reusable new input is
Lemma~\ref{lem:fractional-edge}, followed by the measure change
\eqref{eq:measure-change} and the elementary scalar Lemma~\ref{lem:scalar}.
The existing odd-cycle theorem is used with the same graphon but the biased
probability measure $\nu$.  The graph-specific module awaiting inspection is
\texttt{lean/Taeyoung/Examples/Graph104.lean}.

The method covers exactly one pendant leaf on an odd cycle.  With two or
more leaves, averaging gives a different fractional degree product and the
present scalar transfer does not apply automatically.  It also does not
prove leaf attachment to an arbitrary positive graph: cyclic symmetry is
used essentially in \eqref{eq:agm}.

The chromatic and scalar identities, the fractional H\"older exponents, and
the Atlas identification are audited exactly by
\begin{verbatim}
python codes/verify_odd_cycle_one_leaf.py
\end{verbatim}

%%%%%%%%%%%%%%%%%%%%%%%%%%%%%%%%%%%%%%%%%%%%%%%%%%%%%%%%%%%%%%%%%%%%%%%%%%%%%%%
\section{What the Lean formalization actually proves}
%%%%%%%%%%%%%%%%%%%%%%%%%%%%%%%%%%%%%%%%%%%%%%%%%%%%%%%%%%%%%%%%%%%%%%%%%%%%%%%

The machine-checked development is \texttt{lean/Taeyoung/Methods/OddLeaf/},
ending at \texttt{satisfiesLowerBound\_104}.

\paragraph{Scope: $m=5$ only.}
Theorem~\ref{thm:main} is \emph{not} formalized for general odd $m$.  What is
proved is Corollary~\ref{cor:atlas104}, on all of $p\in[1/2,1]$, and it is the
only member of the family inside the six-vertex catalogue.  The reason is
\eqref{eq:odd-cycle-input}: the project's checked odd-cycle input is the
five-cycle theorem \texttt{OddCycleC5.c5\_homDensity\_bound}, not a general odd
$m$.  The degree-bias construction is nonetheless parametric --- the file
\texttt{OddLeaf/Bias.lean} builds $d^{1/m}$, $M=\int d^{1/m}$, the tilted
measure $\nu$ and the reread kernel for every $m\geq1$ --- while
\texttt{OddLeaf/Holder.lean} and the scalar rescaling are specialized to $m=5$.

\paragraph{The departure: the rescaling carries no derivative.}
Lemma~\ref{lem:scalar} is proved here by differentiating
$\psi_m(s)=\varphi_m(s)s^{-m/2}$, a function with two half-integer powers.  The
formalization forms no half-integer power and differentiates nothing.  Write
\[
 \varphi_5(u)=u^5-u(1-u)^4=u\,h(u),
 \qquad h(u)=(2u-1)(2u^2-2u+1),
\]
which is \texttt{OddLeaf.phi\_eq}, a \texttt{ring} identity.  All of
Lemma~\ref{lem:scalar} then reduces to the polynomial statement
\begin{equation}
 s^3h(p)^2\leq p^3h(s)^2
 \qquad\bigl(\tfrac12\leq p\leq s\bigr),
 \label{eq:lean-oddleaf-mono}
\end{equation}
which is \texttt{OddLeaf.scalar\_mono}.  Substituting $p=\frac12+a^2$ and
$s=\frac12+a^2+b^2$ --- legitimate because the two hypotheses are exactly
$p-\frac12\geq0$ and $s-p\geq0$, so both differences are squares of their own
square roots --- the cleared form $8\bigl(p^3h(s)^2-s^3h(p)^2\bigr)$ becomes a
polynomial in $a,b$ with $36$ terms, every coefficient positive and every
exponent even.  One \texttt{ring} identity and \texttt{positivity} finish it.

\paragraph{Monotonicity of $\varphi_m$ is never used.}
The unnumbered remark after Lemma~\ref{lem:scalar} --- that $\varphi_m'>0$ on
$[1/2,1]$, used in the proof of Theorem~\ref{thm:main} to pass from $s$ to
$s_0$ --- is not formalized.  It is not needed.  Writing $K=M^5$ and letting $s$
be the biased edge density, \texttt{OddLeaf.transfer} proves
\[
 p\,\varphi_5(p)\leq K\,\varphi_5(s)
\]
directly from the three facts $0<K\leq p$, $p^7\leq K^2s^5$ and
\eqref{eq:lean-oddleaf-mono}, by comparing squares:
\[
 \bigl(K\varphi_5(s)\bigr)^2s^5p^5
 =\bigl(K^2s^5\bigr)\bigl(p^5\varphi_5(s)^2\bigr)
 \geq p^7\bigl(s^5\varphi_5(p)^2\bigr)
 =\bigl(p\varphi_5(p)\bigr)^2s^5p^5 .
\]
The intermediate quantity $s_0=p(p/z)^{2/m}$ therefore never appears, and
neither does the hypothesis $s_0\leq1$.  The inequality $p\leq s$, which the
note gets from $z\leq p$ and the definition of $s_0$, is obtained instead from
$p^7\leq K^2s^5\leq p^2s^5$.

\paragraph{Two statements are kept polynomial on purpose.}
Lemma~\ref{lem:fractional-edge} is formalized as
\[
 p^7\leq N^5
 \qquad(\texttt{OddLeaf.pow\_seven\_le\_pow\_rootEdge}),
\]
not as $N\geq p^{7/5}$, so that no real power of $p$ is ever formed; the
$\mathbb{R}_{\geq0}^\infty$ Hölder step itself is exactly the note's, at
exponents $(5/7,1/7,1/7)$, and it reuses the scaffolding of
\texttt{Methods/PathSidorenko.lean} --- the pushforwards \texttt{edgeE},
\texttt{degE}, the two bounds $\int\!\!\int W/d\leq1$, and the null set where
the factorization breaks --- verbatim.  Likewise the concavity bound
$M\leq p^{1/m}$ of the main proof is formalized as $M^m\leq p$
(\texttt{OddLeaf.pow\_rootMean\_le}), by Jensen against the \emph{convex}
$t\mapsto t^m$ rather than by concavity of $r\mapsto r^{1/m}$.  This is why the
exponent is fixed to a reciprocal integer: it keeps both $(d^{1/m})^m=d$ and
$M^m\leq p$ as ordinary $\mathbb{N}$-powers.

\paragraph{The cyclic averaging is a measure-preserving permutation.}
``Cyclic relabelling shows that the same density results when the degree factor
is placed at any $x_i$'' is discharged by
\texttt{Methods/Peeling.integral\_assignment\_perm}, a new general lemma
extracting the measure-preserving half of \texttt{homDensity\_iso} so that it
applies to an integrand which is \emph{not} a graph weight --- here the cycle
weight multiplied by a degree at one coordinate.  It is applied to the rotation
$j\mapsto j+i$, checked to be an automorphism of $C_5$ by \texttt{decide}.  The
arithmetic--geometric mean step \eqref{eq:agm} is Mathlib's
\texttt{Real.geom\_mean\_le\_arith\_mean\_weighted} at the five equal weights
$1/5$, and the measure change \eqref{eq:measure-change} goes through
\texttt{Foundation/TiltTransfer.lean}, the same route as the $d$- and
$\sqrt d$-biases already in the project.

\end{document}
